\documentclass[11pt]{article}
\ifdefined\pdfinfoomitdate\pdfinfoomitdate=1\fi
\ifdefined\pdfsuppressptexinfo\pdfsuppressptexinfo=15\fi
\ifdefined\pdftrailerid\pdftrailerid{}\fi
\usepackage[a4paper,margin=25mm]{geometry}
\usepackage{amsmath,amssymb,amsthm,mathtools,bm}
\usepackage{booktabs,array,microtype}
\usepackage[hidelinks]{hyperref}
\usepackage[nameinlink,noabbrev]{cleveref}
\hypersetup{
  pdftitle={Bayesian Matroid-Union Bounds for Quantum List Discrimination and Exact Adaptive-Parallel Phases},
  pdfauthor={Lluis Eriksson},
  pdfsubject={Support congestion, process compression, and exact adaptive-parallel phases for quantum list discrimination},
  pdfkeywords={quantum state exclusion, list discrimination, matroid union, Rado matroid, channel discrimination, adaptive queries}
}

\newtheorem{theorem}{Theorem}[section]
\newtheorem{proposition}[theorem]{Proposition}
\newtheorem{corollary}[theorem]{Corollary}
\newtheorem{lemma}[theorem]{Lemma}
\theoremstyle{definition}
\newtheorem{definition}[theorem]{Definition}
\theoremstyle{remark}
\newtheorem{remark}[theorem]{Remark}
\newtheorem{example}[theorem]{Example}

\newcommand{\C}{\mathbb C}
\newcommand{\cH}{\mathcal H}
\newcommand{\cI}{\mathcal I}
\newcommand{\cK}{\mathcal K}
\newcommand{\cL}{\mathcal L}
\newcommand{\cM}{\mathcal M}
\newcommand{\cP}{\mathcal P}
\newcommand{\cR}{\mathcal R}
\newcommand{\ran}{\operatorname{ran}}
\newcommand{\rank}{\operatorname{rank}}
\newcommand{\supp}{\operatorname{supp}}
\newcommand{\Tr}{\operatorname{Tr}}
\newcommand{\conv}{\operatorname{conv}}
\newcommand{\ket}[1]{\lvert #1\rangle}
\newcommand{\bra}[1]{\langle #1\rvert}
\newcommand{\vket}[1]{\lvert #1\rangle\!\rangle}
\newcommand{\vbra}[1]{\langle\!\langle #1\rvert}

\title{\textbf{Bayesian Matroid-Union Bounds for Quantum List Discrimination}\\
Support Congestion, Process Compression, and Exact Adaptive--Parallel Phases}
\author{Lluis Eriksson}
\date{August 2026}

\begin{document}
\maketitle

\begin{abstract}
In quantum list discrimination a measurement returns at most $\ell$ candidate
labels and succeeds when the true hypothesis belongs to the returned list.
We derive a hierarchy of upper bounds that retains
the complete labeled dependence pattern of mixed quantum hypotheses and
arbitrary priors.  To the support subspaces of the hypotheses we associate
their Rado independent-transversal matroid $\cR$.  We prove that the vector of
true-label inclusion probabilities of every $\ell$-list measurement belongs
to the independence polytope of the $\ell$-fold matroid union
$\cR^{\vee\ell}$.  Its rank has the closed support-congestion formula
\[
 r_\ell(A)=\min_{C\subseteq A}\bigl(|A\setminus C|
                  +\ell\dim\sum_{i\in C}\ran\rho_i\bigr).
\]
Thus the Bayes list success is at most the maximum prior weight of a set
partitionable into $\ell$ independent support transversals.  The result also
covers soft lists whose fractional rewards have total budget at most $\ell$
per outcome.  A prefix formula gives the bound for arbitrary priors and an exact
rank-saturation audit for equality of any proposed strategy.

For uniform priors the rank formula yields an integer failure certificate
$\Delta_\ell=\max_C(|C|-\ell d(C))$ and the precise list size at which this
support obstruction disappears.  Within any process-tester class satisfying
the stated unnormalised-Choi trace interface, canonical compression transfers
the hierarchy without adding a causal-order assumption.  We exhibit exact
attainment families.  Orthogonal classes
of repeated states attain the arbitrary-prior bound for every $\ell$.
More geometrically, $m$ equiprobable regular-simplex states have exact
one-guess optimum $(m-1)/m$ but perfect two-list success: a root-frame POVM
reports an edge of the simplex and never omits the true vertex.  Finally, an
exact three-qubit-state example shows that a full union matroid does not imply
the existence of a perfect-list measurement; its uniform two-list optimum is
$(2+1/\sqrt3)/3<1$.  The theorem is therefore a
support obstruction with exact attainments, not a general solution of state
exclusion.  The proof
combines standard Rado--Edmonds theory with quantum support inequalities; the
ingredients themselves are not claimed as new.

The process extension becomes exact for a complete input-dependent family
whose channels dephase before preparing a classical output.  For $M=2^h$
equiprobable hypotheses, list size $\ell=2^s$, and $q$ uses, we prove against
arbitrary entangled parallel inputs and adaptive quantum memories that
\[
 P_{\rm par}^{(\ell)}(q)=\min\{1,\ell(q+1)/M\},\qquad
 P_{\rm ad}^{(\ell)}(q)=\min\{1,\ell2^q/M\}.
\]
Thus perfect list identification needs exactly $M/\ell-1$ parallel calls but
only $\log_2(M/\ell)$ adaptive calls.  The fixed parallel transcripts induce
partition matroids whose union caps are attained cell by cell.  This is an
exact list/feedback tradeoff in classical channels embedded quantumly, not a
claim of quantum or indefinite-causal-order advantage.

For a complete unitary-error ensemble and fixed pure probe, we translate the
known optimal approximate dense-coding law into the present process-list
language,
$P_{\rm guess}=(\Tr\sqrt{\rho_R})^2/D$.  It yields the sharp
Schmidt-rank cap $r/D$ and, for independent Weyl histories, exact list optima
for carrier-only serial, unreferenced parallel, and Bell-assisted wirings.
These channel theorems turn process compression into resource-resolved exact
statements while retaining explicit boundaries on prior art and scope.
\end{abstract}

\section{Introduction}

Minimum-error quantum discrimination normally asks a measurement to output
one label.  A different and useful task permits a set-valued decision: the
measurement may return a short list, and it succeeds whenever the true label
is contained in that list.  Taking complements turns this into quantum state
exclusion or elimination.  That subject has a well-developed semidefinite
formulation and includes unambiguous, minimum-error, and worst-case
variants~\cite{bandyopadhyay2014}.  Multi-state elimination has also been
implemented experimentally~\cite{webb2023}.  We do not claim the operational
task as new.

The question addressed here is structural.  Suppose the hypotheses are mixed
and their support subspaces overlap in a complicated labeled pattern.  How
large can the total probability be that the correct label appears in a list
of size at most $\ell$?  A total Hilbert-space dimension records only one
number.  It does not distinguish a dense cluster of dependent, high-prior
labels from independent, low-prior outliers.  In the one-label problem this
loss is repaired by a Rado-matroid Bayes bound~\cite{eriksson2026matroid}.
Here we show that the full list hierarchy is governed not by an ad hoc
multiple of that bound, but by the successive union powers of the same
matroid.

The connection is exact at the polyhedral level.  If $s_i$ is the conditional
probability that label $i$ appears in the reported list, then for each set
$A$ of labels,
\[
  s(A):=\sum_{i\in A}s_i
  \leq \ell\dim\sum_{i\in A}\ran\rho_i.
\]
Combining these support inequalities with $s_i\leq1$ and minimizing over a
subset produces precisely the rank inequalities of the $\ell$-fold union of
the Rado matroid.  Edmonds' union theorem and independence-polytope theorem
then convert the quantum inequality into a maximum-weight independent-set
bound.  This provides, in one statement, arbitrary priors, every finite list
budget, mixed supports, and soft fractional rewards.

Three consequences make the hierarchy operational rather than merely
formal.  First, for equal priors its failure deficit is the integer
\[
 \Delta_\ell=\max_{C}\bigl(|C|-\ell d(C)\bigr),
 \qquad d(C)=\dim\sum_{i\in C}\ran\rho_i.
\]
Second, this matroid support obstruction vanishes at the explicit density
threshold
\[
 \max_C\left\lceil\frac{|C|}{d(C)}\right\rceil.
\]
Third, canonical inverse-square-root
compression carries the argument unchanged to quantum-process testers.  A
fixed tester normalization may describe parallel circuits, adaptive combs,
or more general process-matrix strategies~\cite{chiribella2009,sedlak2016,bavaresco2022}.

Compression by itself is only a converse and does not compare architectures.
We therefore solve one input-dependent channel family exactly.  The unknown
channel measures a query from a binary-prefix alphabet and returns membership
in the corresponding right-child cylinder.  These cylinders are laminar.
For a fixed parallel query tuple, their response fibres form a partition
matroid, so the union bound is attained by reporting the largest-prior labels
in each fibre.  A flag-genie argument proves that coherent or entangled inputs
cannot improve this value, while backward induction reduces every adaptive
quantum strategy to a classical decision tree.  The resulting all-depth phase
diagram separates the linear growth of parallel transcript cells from the
doubling supplied by feedback.  Adaptive advantages for channel
discrimination, including entanglement-breaking examples, are established in
the literature~\cite{harrow2010,bavaresco2021}; our contribution is not the
existence of adaptivity, but the closed list-valued law and its exact
matroid-union realization.

The disappearance threshold of this obstruction is necessary, not sufficient, for perfect list
discrimination.  Geometry inside the support remains relevant.  We make that
limitation exact: three pairwise independent qubit states have a free
two-fold union matroid but cannot be perfectly two-list discriminated because
their Bloch vectors lie in one strict hemisphere.  Conversely, regular
simplices attain a striking sharp transition.  Their one-guess optimum loses
exactly $1/m$, whereas a two-list root-frame measurement succeeds with
certainty for every $m\geq3$.  The channel sections then go beyond static
attainment: one family isolates feedback exactly, while complete unitary-error
ensembles expose how the full Schmidt spectrum and the multitime wiring alter
the achievable list probability.

The paper is organized as follows.  \Cref{sec:model} defines list and soft
reward decisions.  \Cref{sec:matroid} reviews only the Rado and union facts
needed later.  \Cref{sec:main} proves the polytope theorem and the nonuniform
Bayes formula.  \Cref{sec:process} gives the process-tester extension.
\Cref{sec:adaptive} derives the exact parallel--adaptive channel phase.
\Cref{sec:ueb} derives the process-list consequences of the dense-coding
spectrum law and proves the unitary-history wiring laws.
\Cref{sec:deficit} derives the congestion deficit and obstruction threshold.
\Cref{sec:sharp} proves exact attainment families, and
\Cref{sec:counterexample} proves the insufficiency counterexample.  We close
with computational certificates, related work, and explicit limitations.

\section{List decisions and soft rewards}\label{sec:model}

Let $E=\{1,\ldots,M\}$ label density operators $\rho_i$ on a
finite-dimensional Hilbert space $\cK$.  Let $p_i\geq0$ and
$\sum_i p_i=1$.  Zero priors are allowed.

\begin{definition}[Deterministic list measurement]
For an integer $1\leq\ell\leq M$, an $\ell$-list measurement is a POVM
$(M_L)_{L\in\cL}$ indexed by a family
$\cL\subseteq\{L\subseteq E:|L|\leq\ell\}$.  Outcome $L$ reports the labels
in $L$.  Its true-label inclusion probabilities and Bayes list success are
\begin{equation}\label{eq:si}
 s_i=\sum_{L\ni i}\Tr(M_L\rho_i),\qquad
 P_{\rm list}=\sum_i p_i s_i.
\end{equation}
Empty lists may be omitted without changing any statement.
\end{definition}

This convention includes ordinary minimum-error discrimination at $\ell=1$.
At $\ell=M-1$, perfect list success is the usual conclusive exclusion of at
least one label.  Intermediate values quantify multi-label elimination.

The proof needs only a fractional reward budget, not deterministic membership.

\begin{definition}[Soft-list reward rule]
Let $(M_a)_a$ be a POVM.  A soft-list reward rule assigns numbers
$w_{ia}\in[0,1]$ such that for each outcome $a$,
\begin{equation}\label{eq:soft-budget}
 \sum_i w_{ia}\leq\ell.
\end{equation}
Define
\begin{equation}\label{eq:soft}
 s_i=\sum_a w_{ia}\Tr(M_a\rho_i),\qquad
 P=\sum_i p_i s_i.
\end{equation}
\end{definition}

A deterministic list is recovered from $w_{iL}=\bm1_{\{i\in L\}}$.
Fractional $w_{ia}$ allow randomized postprocessing or partial rewards.  The
budget \eqref{eq:soft-budget} is strictly more general than requiring at most
$\ell$ nonzero rewards.  The assumption $w_{ia}\leq1$ is essential for the
singleton bound $s_i\leq1$.

\begin{lemma}[Lists and exclusion]\label{lem:padding}
For a nonnegative inclusion reward, an at-most-$\ell$ list may be padded to
an exactly-$\ell$ list without decreasing success.  Complementing that list
is therefore equivalent to reporting exactly $k=M-\ell$ excluded labels.
In particular, perfect at-most-$\ell$ list discrimination is equivalent to
perfect weak $k$-state exclusion after padding.
\end{lemma}

\begin{proof}
For every outcome list $L$ with $|L|<\ell$, add arbitrary labels from
$E\setminus L$ until its size is $\ell$.  A true label already in $L$ remains
present, and some previously omitted labels may become successful, so the
inclusion reward cannot decrease.  The complement of the padded list has
cardinality $M-\ell$, and the true label lies in the list exactly when it is
not among the excluded labels.  Complementation is reversible.
\end{proof}

This elementary dictionary aligns our convention with $k$-state exclusion
and zero-error list decoding~\cite{stratton2024,dalai2026}.  Stronger exclusion
notions that require exhaustive elimination of every tuple impose additional
conditions and are not identified with the present average-reward task
~\cite{manna2026}.

\section{Rado matroids and their union powers}\label{sec:matroid}

Put $S_i=\ran\rho_i$ and, for $A\subseteq E$, define
\begin{equation}\label{eq:d}
 d(A)=\dim\sum_{i\in A}S_i,\qquad d(\varnothing)=0.
\end{equation}

\begin{definition}[Support Rado matroid]
The support Rado matroid $\cR$ on $E$ declares $I\subseteq E$ independent
when vectors $v_i\in S_i$ can be chosen so that $(v_i)_{i\in I}$ is linearly
independent.
\end{definition}

Rado's independent-transversal theorem~\cite{rado1942} gives the rank formula
\begin{equation}\label{eq:rado-rank}
 r_{\cR}(A)=\min_{C\subseteq A}\bigl(|A\setminus C|+d(C)\bigr).
\end{equation}
Equivalently, $I$ is independent iff $|C|\leq d(C)$ for every $C\subseteq I$.
For rank-one states, $\cR$ is simply the vector matroid of the state vectors.

The union $\cM_1\vee\cM_2$ of two matroids on the same ground set has as
independent sets the unions $I_1\cup I_2$ with $I_j$ independent in
$\cM_j$.  Let $\cR^{\vee\ell}$ denote the union of $\ell$ copies of $\cR$.
Thus $I$ is independent in $\cR^{\vee\ell}$ iff it can be partitioned into
at most $\ell$ $\cR$-independent sets.  Edmonds' matroid-union theorem
~\cite{edmonds1965,edmonds1968} states
\begin{equation}\label{eq:union-standard}
 r_{\cR^{\vee\ell}}(A)
 =\min_{B\subseteq A}\bigl(|A\setminus B|+\ell r_{\cR}(B)\bigr).
\end{equation}

For the present subspace family this formula collapses to one minimization.

\begin{lemma}[Closed union-rank law]\label{lem:union-rank}
For every $A\subseteq E$,
\begin{equation}\label{eq:union-rank}
 r_\ell(A):=r_{\cR^{\vee\ell}}(A)
 =\min_{C\subseteq A}\bigl(|A\setminus C|+\ell d(C)\bigr).
\end{equation}
\end{lemma}

\begin{proof}
Insert \eqref{eq:rado-rank} into \eqref{eq:union-standard}.  The resulting
minimum is over $C\subseteq B\subseteq A$.  Write $B=C\mathbin{\dot\cup}F$.
Then
\begin{align*}
 |A\setminus B|+\ell|B\setminus C|+\ell d(C)
 &=|A\setminus C|+(\ell-1)|F|+\ell d(C).
\end{align*}
For fixed $C$ this is minimized at $F=\varnothing$.  Minimizing over $C$
proves the claim.  The argument also covers $\ell=1$, when every $F$ gives
the same value.
\end{proof}

For a matroid $\cM$ with rank $r$, its independence polytope is
\begin{equation}\label{eq:polytope}
 \cP(\cM)=\{x\in\mathbb R_{
 \geq0}^{E}:x(A)\leq r(A)\text{ for all }A\subseteq E\}.
\end{equation}
Edmonds proved that this polytope is the convex hull of incidence vectors of
independent sets~\cite{edmonds1970}.  Consequently a nonnegative linear
functional is maximized by the matroid greedy algorithm.

\section{The matroid-union list theorem}\label{sec:main}

\begin{theorem}[Support-congestion law]\label{thm:main}
For every density ensemble $(p_i,\rho_i)$, every $\ell\geq1$, and every
soft-list reward rule of budget $\ell$, the vector $s=(s_i)$ in
\eqref{eq:soft} belongs to $\cP(\cR^{\vee\ell})$.  Equivalently, for all
$A\subseteq E$,
\begin{equation}\label{eq:main-subset}
 \sum_{i\in A}s_i
 \leq r_\ell(A)
 =\min_{C\subseteq A}\bigl(|A\setminus C|+\ell d(C)\bigr).
\end{equation}
In particular, the statement holds for deterministic $\ell$-list
measurements.
\end{theorem}

\begin{proof}
Let $P_A$ denote the orthogonal projector onto
$\sum_{i\in A}S_i$.  Since $\rho_i$ is positive with trace one, all its
eigenvalues lie in $[0,1]$.  Therefore
\begin{equation}\label{eq:rho-projector}
 0\leq\rho_i\leq P_A\qquad(i\in A).
\end{equation}
For each outcome $a$, \eqref{eq:rho-projector} and
\eqref{eq:soft-budget} give
\[
 0\leq\sum_{i\in A}w_{ia}\rho_i
 \leq\left(\sum_{i\in A}w_{ia}\right)P_A
 \leq\ell P_A.
\]
It follows that
\begin{align}
 s(A)
 &=\sum_a\Tr\!\left(M_a\sum_{i\in A}w_{ia}\rho_i\right)\notag\\
 &\leq\ell\sum_a\Tr(M_aP_A)
 =\ell\Tr P_A
 =\ell d(A).                                      \label{eq:raw-support}
\end{align}
Also
\[
 0\leq s_i=\Tr\!\left(\rho_i\sum_a w_{ia}M_a\right)\leq1,
\]
because $\sum_a w_{ia}M_a\leq\sum_aM_a=I$.  Hence for arbitrary
$C\subseteq A$,
\[
 s(A)=s(C)+s(A\setminus C)
 \leq\ell d(C)+|A\setminus C|.
\]
Minimizing over $C$ and applying \Cref{lem:union-rank} proves
\eqref{eq:main-subset}.  Nonnegativity of $s$ and the rank-inequality
description \eqref{eq:polytope} finish the proof.
\end{proof}

The raw inequality \eqref{eq:raw-support} alone is not generally a matroid
rank inequality because it can exceed $|A|$.  The singleton constraints are
what perform the exact submodular truncation in \eqref{eq:union-rank}.

\begin{corollary}[Arbitrary-prior Bayes cap]\label{cor:bayes}
Every soft-list rule of budget $\ell$ obeys
\begin{equation}\label{eq:beta}
 P\leq\beta_\ell(p,S)
 :=\max\left\{\sum_{i\in I}p_i:
 I\in\cI(\cR^{\vee\ell})\right\}.
\end{equation}
Thus the cap is the maximum prior mass of a set partitionable into at most
$\ell$ independent support transversals.
\end{corollary}

\begin{proof}
By \Cref{thm:main}, $s\in\cP(\cR^{\vee\ell})$.  The vertices of this
polytope are incidence vectors of independent sets, so maximizing $p\cdot s$
over it gives \eqref{eq:beta}.
\end{proof}

The cap has a useful formula that avoids choosing a maximizing independent
set explicitly.  Order the labels so that
$p_{(1)}\geq\cdots\geq p_{(M)}$ and put
$E_j=\{(1),\ldots,(j)\}$, $E_0=\varnothing$, and $p_{(M+1)}=0$.

\begin{proposition}[Prefix formula and equality audit]\label{prop:prefix}
One has
\begin{equation}\label{eq:prefix}
 \beta_\ell(p,S)=
 \sum_{j=1}^{M}\bigl(p_{(j)}-p_{(j+1)}\bigr)r_\ell(E_j).
\end{equation}
For a fixed soft-list rule with inclusion vector $s$, equality
$p\cdot s=\beta_\ell$ holds iff
\begin{equation}\label{eq:prefix-eq}
 s(E_j)=r_\ell(E_j)
 \quad\text{whenever }p_{(j)}>p_{(j+1)}.
\end{equation}
The condition is independent of how labels are ordered inside a prior tie.
\end{proposition}

\begin{proof}
The matroid greedy theorem selects label $(j)$ precisely when
$r_\ell(E_j)-r_\ell(E_{j-1})=1$.  Summation by parts gives
\eqref{eq:prefix}.  The same identity for an arbitrary vector is
\[
 p\cdot s=\sum_{j=1}^{M}
 (p_{(j)}-p_{(j+1)})s(E_j).
\]
Every coefficient and every deficit
$r_\ell(E_j)-s(E_j)$ is nonnegative.  Their weighted sum vanishes exactly at
the strict prior breakpoints, proving \eqref{eq:prefix-eq}.  At a tie the
coefficient is zero, so rearranging tied labels has no effect.
\end{proof}

\begin{remark}[Hierarchy]
The matroids satisfy
$\cR\preceq\cR^{\vee2}\preceq\cdots$, so
$\beta_1\leq\beta_2\leq\cdots\leq1$.  The case $\ell=1$ is the previously
derived one-label Rado bound~\cite{eriksson2026matroid}.  The new content is
the union hierarchy, its soft-decision meaning, and the exact obstruction threshold and
attainment results below.
\end{remark}

\section{General quantum-process testers}\label{sec:process}

We use unnormalised Choi operators throughout this section.  This avoids an
input-dimension factor in the tester normalization.

Let $J_i\geq0$ be the Choi operators of a finite family of deterministic
processes.  A tester with list-valued outcomes consists of positive operators
$T_L$ whose sum $T=\sum_LT_L$ is a valid deterministic tester normalization.
For every candidate process in the stated tester interface,
\begin{equation}\label{eq:tester-norm}
 \Tr(TJ_i)=1,
\end{equation}
and the conditional probability of outcome $L$ is $\Tr(T_LJ_i)$.  The
linear constraints selecting parallel, sequential, or general testers affect
the admissible set of $T$, but not the following compression.

\begin{lemma}[Canonical compression]\label{lem:compression}
Let $\cK=\supp T$ and let $T^{-1/2}$ be the Moore--Penrose inverse on
$\cK$.  Define
\begin{equation}\label{eq:compression}
 \rho_i=T^{1/2}J_iT^{1/2},\qquad
 M_L=T^{-1/2}T_LT^{-1/2}\big|_{\cK}.
\end{equation}
Then $(\rho_i)$ are density operators on $\cK$, $(M_L)$ is a POVM on
$\cK$, and
\begin{equation}\label{eq:prob-preserve}
 \Tr(M_L\rho_i)=\Tr(T_LJ_i)
\end{equation}
for all $i,L$.
\end{lemma}

\begin{proof}
Positivity is immediate and \eqref{eq:tester-norm} gives
$\Tr\rho_i=1$.  Since $0\leq T_L\leq T$, the support of every $T_L$ is
contained in $\cK$, and
\[
 \sum_LM_L=T^{-1/2}TT^{-1/2}=I_{\cK}.
\]
Using the support projection $P_{\cK}=T^{-1/2}T^{1/2}$, cyclicity of trace,
and $T_L=P_{\cK}T_LP_{\cK}$ gives
\begin{align*}
 \Tr(M_L\rho_i)
 &=\Tr(T^{-1/2}T_LT^{-1/2}T^{1/2}J_iT^{1/2})\\
 &=\Tr(P_{\cK}T_LP_{\cK}J_i)=\Tr(T_LJ_i).
\end{align*}
\end{proof}

\begin{theorem}[Process-list matroid bound]\label{thm:process}
Fix any admissible tester normalization $T$ satisfying
\eqref{eq:tester-norm}.  Let
$S_i^T=\ran(T^{1/2}J_iT^{1/2})$.  Every $\ell$-list tester with this
normalization obeys \Cref{thm:main,cor:bayes,prop:prefix} for the Rado matroid
of $(S_i^T)$.

Let $S_i^0=\ran J_i$ and
$d_0(A)=\dim\sum_{i\in A}S_i^0$.  Then a tester-independent bound is obtained
by replacing $d_T$ by $d_0$ in \eqref{eq:union-rank}.  In particular, it
holds after optimizing over all admissible general testers.
\end{theorem}

\begin{proof}
The fixed-$T$ statement follows from \Cref{lem:compression,thm:main}.
Moreover,
\[
 \ran(T^{1/2}J_iT^{1/2})\subseteq T^{1/2}\ran J_i,
\]
so a common linear map can only decrease the dimension of every support sum:
$d_T(A)\leq d_0(A)$.  Formula \eqref{eq:union-rank} is monotone in $d$, and
so are the prefix and maximum-weight bounds.  The resulting cap does not
depend on $T$ and survives the outer optimization.
\end{proof}

For rank-one Choi operators $J_i=\vket{U_i}\!\vbra{U_i}$, the oracle matroid
is the vector matroid of the Choi vectors.  The theorem therefore bounds list
discrimination of unitary channels by the union powers of their labeled
linear-dependence matroid.  We make no claim that the tester-independent cap
is always attainable, or that indefinite causal order is irrelevant to the
actual optimum.

For clarity, equality in the effective and original-support caps are different
claims.  For fixed $T$, \Cref{prop:prefix} characterizes equality in the cap
built from $S_i^T$.  Equality in the coarser cap built from $S_i^0$ additionally
requires that compression cause no weighted prefix-rank loss.  We never infer
the latter from saturation of the former.

\subsection{A robust core-and-tail certificate}

Exact support can change discontinuously under a small full-rank perturbation.
The following secondary result retains a low-rank union-matroid core and pays
explicitly for the discarded tail.  It extends the one-label robust
certificate of Ref.~\cite{eriksson2026matroid} to every list budget.

\begin{corollary}[Robust process cores]\label{cor:robust}
Let $0\leq L_i\leq J_i$ and $R_i=J_i-L_i$.  For a fixed deterministic tester
normalization $T$, put
\[
 \lambda_i=T^{1/2}L_iT^{1/2},\qquad S_i^L=\ran\lambda_i,
\]
and form the union-matroid cap $\beta_\ell(T,L,p)$ from the subspaces
$(S_i^L)$.  Every deterministic $\ell$-list tester with normalization $T$
obeys
\begin{equation}\label{eq:robust-fixed}
 P_{\rm list}\leq
 \min\left\{1,\ \beta_\ell(T,L,p)+\sum_i p_i\Tr(TR_i)\right\}.
\end{equation}
Consequently, for an architecture class $\mathfrak T$ of deterministic
normalizations,
\begin{equation}\label{eq:robust-universal}
 P_{\rm list}^{\mathfrak T}\leq\min\left\{1,
 \sup_{T\in\mathfrak T}\beta_\ell(T,L,p)
 +\sum_i p_i\sup_{T\in\mathfrak T}\Tr(TR_i)\right\}.
\end{equation}
The first supremum may be replaced by the coarser cap computed from the
uncompressed core supports $\ran L_i$.
\end{corollary}

\begin{proof}
Decompose $s_i=s_i^L+s_i^R$ using $J_i=L_i+R_i$.  The compressed core obeys
$0\leq\lambda_i\leq\rho_i$, where $\rho_i$ is the density operator from
\Cref{lem:compression}.  Thus $\lambda_i$ is bounded by the projector onto
any core support sum containing it.  The proof of \Cref{thm:main} applies to
the core contributions even though $\Tr\lambda_i$ need not be one, and gives
$\sum_i p_i s_i^L\leq\beta_\ell(T,L,p)$.

For the tail,
\[
 \sum_{L\ni i}T_L\leq\sum_LT_L=T,
\]
so positivity gives $s_i^R\leq\Tr(TR_i)$.  Summing with the priors and capping
a success probability by one proves \eqref{eq:robust-fixed}.  Taking the
architecture supremum and bounding each tail term separately proves
\eqref{eq:robust-universal}.  Common congruence can only lower core
support-sum ranks, proving the final statement.
\end{proof}

\begin{remark}
The tail term in \eqref{eq:robust-universal} deliberately contains separate
suprema.  It does not assert that one tester normalization simultaneously
maximizes all tails.  For ordinary state discrimination, $T=I$ and the
penalty reduces to $\sum_i p_i\Tr R_i$.
\end{remark}

\section{Input-dependent channels and an exact architecture phase}
\label{sec:adaptive}

The tester theorem above bounds each architecture after its deterministic
normalization has been chosen.  It does not normally optimize that choice or
compare architectures.  We now give a channel family for which both tasks
close exactly, including arbitrary entangled parallel probes and arbitrary
adaptive quantum memories.

Let $\mathsf X$ and $\mathsf Y$ be finite alphabets.  Hypothesis $i$ specifies
a deterministic function $f_i:\mathsf X\to\mathsf Y$ and the
entanglement-breaking channel
\begin{equation}\label{eq:classical-channel}
 \Phi_i(\tau)=\sum_{x\in\mathsf X}
  \bra{x}\tau\ket{x}\,
  \ket{f_i(x)}\!\bra{f_i(x)}.
\end{equation}
The same unknown channel may be called repeatedly.  A parallel strategy uses
$q$ copies simultaneously and may prepare an arbitrary state entangled across
the $q$ inputs and a reference.  An adaptive strategy may retain an arbitrary
quantum memory and choose each later input after the earlier classical
outputs.  No indefinite-causal-order strategy is included in the comparison.

For a deterministic query tuple $\bm x=(x_1,\ldots,x_q)$, set
\[
 c_{\bm x}(i)=(f_i(x_1),\ldots,f_i(x_q)),\qquad
 C_z(\bm x)=\{i:c_{\bm x}(i)=z\}.
\]
If $w=(w_i)_{i\in E}$ is a nonnegative weight vector, write
\begin{equation}\label{eq:cell-value}
 B_\ell(\bm x;w)=
 \sum_z\ \sum_{j=1}^{\min\{\ell,|C_z(\bm x)|\}}
 w^{\downarrow}_{C_z(\bm x),j},
\end{equation}
where $w^{\downarrow}_{C,j}$ is the $j$th largest weight in $C$.

\begin{theorem}[Exact reduction of classical channels]
\label{thm:classical-reduction}
For the channels in \eqref{eq:classical-channel} and arbitrary priors $p$,
the optimum of every $q$-use parallel quantum strategy is
\begin{equation}\label{eq:parallel-cell}
 P_{\rm par}^{(\ell)}(q)=
 \max_{\bm x\in\mathsf X^q}B_\ell(\bm x;p).
\end{equation}
The adaptive optimum is the exact Bellman value
\begin{align}
 V_0(w)&=\sum_{j=1}^{\min\{\ell,|\supp w|\}}w_j^\downarrow,
 \label{eq:bellman0}\\
 V_q(w)&=\max_{x\in\mathsf X}\sum_{y\in\mathsf Y}
 V_{q-1}\bigl((w_i\mathbf1_{\{f_i(x)=y\}})_i\bigr),
 \qquad P_{\rm ad}^{(\ell)}(q)=V_q(p).
 \label{eq:bellman}
\end{align}
For fixed $\bm x$, the effective supports form the partition matroid
\begin{equation}\label{eq:partition-matroid}
 \bigoplus_z U_{1,|C_z(\bm x)|}.
\end{equation}
Its $\ell$-fold union bound is exactly \eqref{eq:cell-value} and is attained
by reporting the $\ell$ largest-prior labels in each response cell.
\end{theorem}

\begin{proof}
For deterministic basis inputs, the output word is the orthogonal classical
state $\ket{c_{\bm x}(i)}$.  Labels in one cell are identical and different
cells are perfectly distinguishable.  The cellwise decision just described
is therefore optimal.  Its support matroid is
\eqref{eq:partition-matroid}; taking $\ell$ union copies changes every class
capacity from one to $\min\{\ell,|C_z|\}$, and matroid greedy optimization
gives \eqref{eq:cell-value}.

For a general parallel input on $\mathsf X^{\otimes q}\otimes R$, each call
first measures its input in the fixed basis.  Give the decoder a genie flag
containing the resulting tuple $\bm x$.  Conditional on this flag, the
reference operator is independent of $i$ and all label dependence lies in
$c_{\bm x}(i)$.  Revealing a flag cannot lower success, so the original
strategy is bounded by a convex combination of the values
$B_\ell(\bm x;p)$ and hence by their maximum.  A deterministic basis tuple
attains that maximum, proving \eqref{eq:parallel-cell}.

The same flag may be revealed after every call of a sequential strategy.
Conditional on the complete prior input--output transcript, induction shows
that the retained quantum state is fixed by the strategy and independent of
the label inside the surviving response cell.  Thus every quantum strategy
is bounded by a randomized classical query tree with the same transcript.
At each node the value is affine in the randomized next query, so an extreme
deterministic query is optimal.  The terminal reward is the sum of the
$\ell$ largest surviving weights, and conditioning on the next output gives
\eqref{eq:bellman0}--\eqref{eq:bellman}.  Conversely, every deterministic
query tree is an admissible adaptive quantum strategy, so the bound is exact.
\end{proof}

We now choose a family on which the recursion and the parallel cell problem
both have closed forms.  Fix $h\geq1$ and label $M=2^h$ channels by the
binary leaves $u\in\{0,1\}^h$.  The input alphabet consists of all prefixes
$v$ with $|v|<h$, and
\begin{equation}\label{eq:prefix-channel}
 f_u(v)=\mathbf1_{\{v1\text{ is a prefix of }u\}}.
\end{equation}
Thus query $v$ asks membership in the right child below $v$.  The queried
cylinders $A_v=\{u:f_u(v)=1\}$ are laminar: any two are nested or disjoint.

\begin{theorem}[Complete parallel--adaptive list phase]
\label{thm:laminar-phase}
Let the channels \eqref{eq:prefix-channel} have uniform priors, let
$\ell=2^s$ with $0\leq s\leq h$, and let $q\geq0$.  Then
\begin{equation}\label{eq:architecture-phase}
 P_{\rm par}^{(\ell)}(q)=
 \min\left\{1,\frac{\ell(q+1)}{2^h}\right\},\qquad
 P_{\rm ad}^{(\ell)}(q)=
 \min\left\{1,\frac{\ell2^q}{2^h}\right\}.
\end{equation}
Consequently the exact numbers of calls required for perfect list success are
\begin{equation}\label{eq:query-thresholds}
 q_{\rm par}^*=2^{h-s}-1=\frac{M}{\ell}-1,\qquad
 q_{\rm ad}^*=h-s=\log_2\frac{M}{\ell}.
\end{equation}
In particular, one-label identification needs $2^h-1$ parallel calls but
only $h$ adaptive calls.
\end{theorem}

\begin{table}[ht]
\centering
\caption{Exact resource phase for the laminar channel family, with
$M=2^h$ and $\ell=2^s$.  ``Cells'' denotes the maximum number of classical
response cells before saturation.}
\label{tab:adaptive-phase}
\begin{tabular}{@{}lccc@{}}
\toprule
Architecture & cells after $q$ calls & list success & calls for success one\\
\midrule
parallel, no feedback & $q+1$ & $\min\{1,\ell(q+1)/M\}$ & $M/\ell-1$\\
adaptive feed-forward & $2^q$ & $\min\{1,\ell2^q/M\}$ & $\log_2(M/\ell)$\\
\bottomrule
\end{tabular}
\end{table}

\begin{proof}
Any $q$ fixed queries select $q$ sets from a laminar family.  Such sets have
at most $q+1$ nonempty membership cells: insert them in an order compatible
with inclusion; a new set lies inside one old cell and can split only that
cell.  For uniform priors, \Cref{thm:classical-reduction} assigns at most
$\ell/2^h$ success to each cell.  This proves the parallel upper bound.

Put $d=h-s$.  Query the prefixes of lengths below $d$ in breadth-first order.
Each new right-child cylinder splits one current prefix block into two, and
until all $2^d-1$ queries have been used every resulting block contains at
least $2^s=\ell$ leaves.  After $q\leq2^d-1$ queries there are therefore
exactly $q+1$ cells, each contributing $\ell/2^h$.  At $q=2^d-1$ the cells
are the $2^d$ length-$d$ prefix blocks, each of size $\ell$; later calls can
be ignored.  This attains the first formula.

A binary adaptive tree of depth $q$ has at most $2^q$ transcripts, and each
terminal list covers at most $\ell$ labels.  Hence its uniform success is at
most the second expression in \eqref{eq:architecture-phase}.  To attain it,
start at the empty prefix and query the currently certified prefix $v$.
Output one selects child $v1$, while output zero selects $v0$.  After $q$
rounds the transcript identifies the first $q$ bits and the surviving cell
has size $2^{h-q}$.  At $q=d$ this size is $\ell$, proving both the second
formula and \eqref{eq:query-thresholds}.
\end{proof}

The depth-two member has four channels with function tables
\[
 000,\qquad010,\qquad100,\qquad101.
\]
With two calls, \Cref{thm:laminar-phase} gives
\begin{equation}\label{eq:feedback-list-tradeoff}
 P_{\rm par}^{(1)}=\frac34,\qquad
 P_{\rm par}^{(2)}=1,\qquad
 P_{\rm ad}^{(1)}=1.
\end{equation}
Thus one feedback bit and one extra reported label close the same parallel
one-guess deficit.  The example is entirely classical after the input
measurement; it proves neither a quantum-memory advantage nor an
indefinite-causal-order advantage.  Binary search, membership-query learning,
and adaptive channel advantages are classical or established mechanisms
\cite{angluin1988,harrow2010}.  The result here is the exact list-valued phase
against the full parallel and adaptive quantum strategy classes specified
above, together with its attained union-matroid interpretation.

\section{Entanglement spectra and multitime unitary histories}
\label{sec:ueb}

The preceding family isolates feedback but is classical after dephasing.  We
next give an input-sensitive quantum family in which the complete Schmidt
spectrum, rather than only support rank, determines the optimum.  The result
also separates three precisely delimited multitime wirings.

Let $\dim A=D$ and let $(U_g)_{g=1}^{D^2}$ be a complete unitary error basis,
so
\begin{equation}\label{eq:ueb}
 \Tr(U_g^*U_{g'})=D\delta_{gg'}.
\end{equation}
For a pure probe $\psi\in A\otimes R$, put
$\rho_R=\Tr_A\ket{\psi}\!\bra{\psi}$ and
$\psi_g=(U_g\otimes I)\psi$.

\begin{theorem}[Dense-coding spectrum law in process form]
\label{thm:ueb-spectrum}
For uniform discrimination of the $D^2$ channels
$\Phi_g=\operatorname{Ad}_{U_g}$ with the fixed probe $\psi$,
\begin{equation}\label{eq:ueb-success}
 P_{\rm guess}(\psi)=\frac{(\Tr\sqrt{\rho_R})^2}{D}.
\end{equation}
Consequently, among probes of Schmidt rank at most $r\leq D$,
\begin{equation}\label{eq:schmidt-rank-cap}
 \max_{\operatorname{SR}(\psi)\leq r}P_{\rm guess}(\psi)=\frac rD,
\end{equation}
with equality exactly for a flat nonzero Schmidt spectrum.  Every
$\ell$-list measurement for the same fixed probe satisfies
\begin{equation}\label{eq:ueb-list-cap}
 P_{\rm list}^{(\ell)}(\psi)\leq
 \min\left\{1,\frac{\ell(\Tr\sqrt{\rho_R})^2}{D}\right\}.
\end{equation}
\end{theorem}

The one-guess identity \eqref{eq:ueb-success}, including optimality for an
arbitrary pure entangled resource, is the approximate dense-coding result of
Feng, Duan, and Ji~\cite{feng2006}.  We include the short frame/dual proof to
fix conventions and to make the list and multitime consequences below
self-contained; we do not claim the identity itself as new.

\begin{proof}
Hilbert--Schmidt completeness gives the depolarizing identity
\[
 \sum_gU_g XU_g^*=D\Tr(X)I_A.
\]
Therefore, on $A\otimes\supp\rho_R$,
\begin{equation}\label{eq:ueb-frame}
 S:=\sum_g\ket{\psi_g}\!\bra{\psi_g}
   =D I_A\otimes\rho_R.
\end{equation}
The square-root measurement
$M_g=S^{-1/2}\ket{\psi_g}\!\bra{\psi_g}S^{-1/2}$ sums to the identity on
this support.  Moreover
\[
 c:=\bra{\psi_g}S^{-1/2}\ket{\psi_g}
   =\frac{\Tr\sqrt{\rho_R}}{\sqrt D}
\]
is independent of $g$, so its success is $c^2$.

For the dual minimum-error SDP, set
$Y=(c/D^2)S^{1/2}$.  The rank-one domination criterion gives
$Y\succeq\ket{\psi_g}\!\bra{\psi_g}/D^2$, because
\[
 \bra{\psi_g}(cS^{1/2})^{-1}\ket{\psi_g}=1.
\]
Also $\Tr Y=(\Tr\sqrt{\rho_R})^2/D$, proving optimality.  If the nonzero
Schmidt coefficients are $\lambda_1,\ldots,\lambda_r$, Cauchy--Schwarz gives
$(\sum_j\sqrt{\lambda_j})^2\leq r$, with equality exactly when they are
all $1/r$.

Finally, for a list POVM $(M_L)$ define
$N_g=\ell^{-1}\sum_{L\ni g}M_L$.  Then
$\sum_gN_g\preceq I$; completing this sub-POVM to a one-label POVM shows
$P_{\rm guess}\geq P_{\rm list}^{(\ell)}/\ell$.  Combining with
\eqref{eq:ueb-success} and the trivial cap one proves
\eqref{eq:ueb-list-cap}.
\end{proof}

The formula is stable under independent time slots.  A tensor product of
complete unitary error bases is a complete basis on dimension
$D=\prod_td_t$, so \Cref{thm:ueb-spectrum} already allows inputs entangled
across different times.  In particular, without a retained reference every
pure parallel input has success $1/D$, whereas a global Schmidt-rank budget
$r$ gives exact optimum $\min\{r,D\}/D$.

For a list-valued wiring comparison, let each of $n$ time slots carry an
independent uniformly labeled Weyl operator $X^{a_t}Z^{b_t}$ on $\C^d$.
There are $M=d^{2n}$ history hypotheses.  We compare:
\begin{itemize}
 \item $\mathrm{SER}_0$: one $d$-dimensional carrier is reused, with known
       interleaving unitaries but no side memory, intermediate measurement,
       or retained transcript;
 \item $\mathrm{PAR}_0$: $n$ fresh inputs are used in parallel and may be
       entangled with each other, but no reference is retained;
 \item $\mathrm{PAR}_{\rm Bell}$: every input has a retained
       $d$-dimensional Bell reference.
\end{itemize}

\begin{theorem}[Exact multitime wiring trichotomy]
\label{thm:weyl-wiring}
For the uniform Weyl histories above,
\begin{equation}\label{eq:weyl-trichotomy}
 \begin{aligned}
 P_{\mathrm{SER}_0}^{(\ell)}
   &=\min\{1,\ell/d^{2n-1}\},\\
 P_{\mathrm{PAR}_0}^{(\ell)}
   &=\min\{1,\ell/d^n\},\\
 P_{\mathrm{PAR}_{\rm Bell}}^{(\ell)}&=1.
 \end{aligned}
\end{equation}
Before saturation, fresh parallel inputs improve on a carrier-only serial
wiring by the exact factor $d^{n-1}$.
\end{theorem}

\begin{table}[ht]
\centering
\caption{Exact list values for $M=d^{2n}$ independent Weyl histories.  The
resource restrictions are part of the architecture definitions preceding
\Cref{thm:weyl-wiring}.}
\label{tab:weyl-wiring}
\begin{tabular}{@{}lcc@{}}
\toprule
Architecture & terminal support dimension & optimal list success\\
\midrule
$\mathrm{SER}_0$ & $d$ & $\min\{1,\ell/d^{2n-1}\}$\\
$\mathrm{PAR}_0$ & $d^n$ & $\min\{1,\ell/d^n\}$\\
$\mathrm{PAR}_{\rm Bell}$ & $d^{2n}$ orthogonal labels & $1$\\
\bottomrule
\end{tabular}
\end{table}

\begin{proof}
A uniform ensemble of $M$ pure states in dimension $q_0$ obeys
$P_{\rm list}^{(\ell)}\leq\min\{1,\ell q_0/M\}$ by
\Cref{thm:main}.  Every $\mathrm{SER}_0$ output occupies only the terminal
$d$-dimensional carrier.  With no interleaving, Weyl multiplication produces
the net label
$(A,B)=(\sum_ta_t,\sum_tb_t)$ modulo $d$; input $\ket0$ and a computational
measurement reveal $A$.  Each value leaves exactly $d^{2n-1}$ histories, so
reporting any $\ell$ of them attains the first formula.

Every $\mathrm{PAR}_0$ output occupies dimension $d^n$, even when its inputs
are entangled across slots.  The product input $\ket0^{\otimes n}$ reveals
the full shift vector $(a_1,\ldots,a_n)$ and leaves exactly $d^n$ phase
histories, attaining the second formula.  Local Bell inputs turn the Weyl
outputs into an orthonormal product Bell basis and prove the third.
\end{proof}

The support-rank upper bound in \eqref{eq:ueb-list-cap} need not be attained
for an intermediate entanglement spectrum.  For example, take $D=4$ and the
rank-two probe $(\ket{00}+\ket{11})/\sqrt2$ for the $16$ Weyl channels.  The
outputs split into four orthogonal shift sectors, each containing the phase
states
\[
 \phi_b=(\ket0+\mathrm{i}^b\ket1)/\sqrt2,
 \qquad b=0,1,2,3.
\]
Adjacent-pair reward operators have largest eigenvalue $1+1/\sqrt2$ and
their top eigenvectors form a tight frame; the covariant dual is
$(1+1/\sqrt2)I/4$.  Pinching over the shift sectors therefore proves the exact
fixed-probe value
\begin{equation}\label{eq:rank-two-list}
 P_{\rm list}^{(2)}=\frac{2+\sqrt2}{4}<1,
\end{equation}
although the coarse right-hand side of \eqref{eq:ueb-list-cap} equals one.
This is a fixed-probe statement, not a resource-optimized impossibility
theorem.

Unitary-error bases, dense coding, group-covariant discrimination, and
parallel treatment of independent unitary channels are established
subjects~\cite{feng2006,sacchi2005,wu2006,chiribella2008,zhuang2020,bavaresco2022}.
We claim neither those ingredients, the one-guess spectrum formula, nor the
dimension bound as new.  The contribution of this section is narrower: the
list-cap embedding, the exact multitime wiring trichotomy, and the explicit
intermediate-spectrum failure of the coarse list cap.

\section{Uniform congestion, deficits, and obstruction thresholds}\label{sec:deficit}

For equal priors, \Cref{cor:bayes} depends only on the rank of the full ground
set.

\begin{corollary}[Exact support deficit]\label{cor:deficit}
For $p_i=1/M$,
\begin{equation}\label{eq:deficit}
 P_{\rm list}\leq\frac{r_\ell(E)}{M}
 =1-\frac{\Delta_\ell}{M},\qquad
 \Delta_\ell=\max_{C\subseteq E}\bigl(|C|-\ell d(C)\bigr).
\end{equation}
The maximum includes $C=\varnothing$, so $\Delta_\ell\geq0$.
\end{corollary}

\begin{proof}
Taking $A=E$ in \eqref{eq:union-rank},
\[
 r_\ell(E)=\min_C(M-|C|+\ell d(C))
 =M-\max_C(|C|-\ell d(C)).
\]
Divide by $M$.
\end{proof}

The integer $\Delta_\ell$ counts the rank deficit of the union matroid, not a
continuous geometric error.  It gives a fail-closed certificate: if any
label cluster has more than $\ell$ labels per available support dimension,
uniform perfect-list discrimination is impossible.

\begin{corollary}[Disappearance threshold of the matroid obstruction]\label{cor:threshold}
The support bound permits $P_{\rm list}=1$ exactly when
\begin{equation}\label{eq:hall-ell}
 |C|\leq\ell d(C)\qquad\text{for all }C\subseteq E.
\end{equation}
The smallest list budget for which this happens is
\begin{equation}\label{eq:ell-support}
 \ell_{\rm supp}
 =\max_{\varnothing\ne C\subseteq E}
   \left\lceil\frac{|C|}{d(C)}\right\rceil.
\end{equation}
Equivalently, $E$ can then be partitioned into $\ell_{\rm supp}$ independent
sets of the support Rado matroid.
\end{corollary}

\begin{proof}
The cap equals one iff $r_\ell(E)=M$, equivalently iff
$\Delta_\ell=0$.  This is exactly \eqref{eq:hall-ell}; minimizing the integer
$\ell$ gives \eqref{eq:ell-support}.  The final statement is Edmonds'
matroid-partition criterion applied to $\cR$, or directly the condition that
$E$ be independent in $\cR^{\vee\ell}$.
\end{proof}

The adjective ``support'' in $\ell_{\rm supp}$ is essential.  It is a lower
bound on the list budget required for perfect physical success, but equality
need not hold; \Cref{sec:counterexample} gives an exact failure.

It is also not the strongest known feasibility obstruction.  The projector
test for conclusive $k$-state exclusion of Stratton, Hsieh, and
Skrzypczyk~\cite{stratton2024}, applied after \Cref{lem:padding} and then to
each restricted label set $C$, requires
\begin{equation}\label{eq:projector-test}
 \sum_{i\in C}\Pi_i\leq\ell P_C,
\end{equation}
where $\Pi_i$ projects onto $S_i$ and $P_C$ onto
$\sum_{i\in C}S_i$.  For pure states, taking the trace recovers
$|C|\leq\ell d(C)$.  For mixed states, and even geometrically for pure states,
the operator inequality can be strictly stronger.  Thus
\eqref{eq:ell-support} is precisely the disappearance threshold of the
\emph{matroid} obstruction, not a new feasibility characterization.

\begin{proposition}[Strict improvement over total dimension]\label{prop:strict-flat}
Let five distinct pure states lie in a common two-dimensional subspace and
let a sixth pure state span an orthogonal line.  Give the five coplanar labels
prior $a=19/100$ each and the orthogonal label prior $b=5/100$.  For list
budget $\ell=2$, the total-dimension relaxation is one, whereas the
matroid-union cap is
\begin{equation}\label{eq:strict-flat}
 \beta_2=4a+b=\frac{81}{100}.
\end{equation}
\end{proposition}

\begin{proof}
Choose, for example,
$\phi_j=(e_0+j e_1)/\sqrt{1+j^2}$ for $j=0,\ldots,4$, and
$\phi_5=e_2$.  The vector matroid is
$U_{2,5}\oplus U_{1,1}$.  Its two-fold union is
$U_{4,5}\oplus U_{1,1}$, so a maximum-weight independent set retains four
coplanar labels and the coloop, proving \eqref{eq:strict-flat}.  The total
dimension is $D=3$ and $\ell D=6=M$; hence the relaxation that keeps only
$D$ and the singleton bounds gives one.  No attainability claim is made for
the value $81/100$.
\end{proof}

\section{Sharp families}\label{sec:sharp}

We first show that the arbitrary-prior theorem is exactly attainable for a
complete family of partition matroids.

\begin{proposition}[Orthogonal parallel classes]\label{prop:parallel}
Let $(e_a)_{a=1}^q$ be orthonormal.  Partition $E$ into nonempty classes
$E_a$, and set $\rho_i=\ket{e_a}\!\bra{e_a}$ for $i\in E_a$.  For every
prior and list budget $\ell$,
\begin{equation}\label{eq:parallel-opt}
 P_{\rm list}^{\rm opt}
 =\sum_{a=1}^q\ 
   \sum_{i\in\operatorname{Top}_{\min(\ell,|E_a|)}(E_a)}p_i.
\end{equation}
This equals the matroid-union cap.
\end{proposition}

\begin{proof}
The support matroid is the partition matroid with capacity one in each class;
its $\ell$-fold union has capacity $\ell$ in each class.  The
maximum-weight independent set therefore selects the $\min(\ell,|E_a|)$
largest priors in each class, proving the upper bound.  Measure the orthogonal
projectors $\ket{e_a}\!\bra{e_a}$ and, on outcome $a$, report precisely those
top-prior labels.  This attains \eqref{eq:parallel-opt}.
\end{proof}

Repeated states make \Cref{prop:parallel} combinatorially transparent.  The
next family uses only distinct, pairwise nonorthogonal pure states.

\subsection{The regular-simplex transition}

Let $m\geq3$ and let unit vectors
$\phi_1,\ldots,\phi_m\in\C^{m-1}$ satisfy
\begin{equation}\label{eq:simplex}
 \langle\phi_i,\phi_j\rangle=-\frac1{m-1}\quad(i\ne j).
\end{equation}
Their Gram matrix has one zero eigenvalue and eigenvalue $m/(m-1)$ with
multiplicity $m-1$.  Consequently
\begin{equation}\label{eq:simplex-frame}
 \sum_{i=1}^m\phi_i=0,\qquad
 \sum_{i=1}^m\ket{\phi_i}\!\bra{\phi_i}
 =\frac{m}{m-1}I.
\end{equation}
Every proper subset is linearly independent, so their vector matroid is the
uniform matroid $U_{m-1,m}$.

\begin{theorem}[Exact one-to-two-list jump]\label{thm:simplex}
For the equiprobable regular-simplex ensemble,
\begin{equation}\label{eq:simplex-opt}
 P_{\rm opt}^{(1)}=\frac{m-1}{m},\qquad
 P_{\rm opt}^{(2)}=1.
\end{equation}
An optimal one-guess POVM is
\begin{equation}\label{eq:one-povm}
 N_i=\frac{m-1}{m}\ket{\phi_i}\!\bra{\phi_i}.
\end{equation}
An optimal two-list POVM is indexed by unordered pairs $i<j$:
\begin{equation}\label{eq:edge-povm}
 \psi_{ij}=\frac{\phi_i-\phi_j}{\sqrt{2m/(m-1)}},\qquad
 M_{ij}=\frac2m\ket{\psi_{ij}}\!\bra{\psi_{ij}},
\end{equation}
and outcome $(i,j)$ reports $\{i,j\}$.
\end{theorem}

\begin{proof}
For one guess, the matroid cap is
$r_1(E)/m=(m-1)/m$.  Equations
\eqref{eq:simplex-frame} and \eqref{eq:one-povm} show that $(N_i)$ is a POVM,
and its conditional correct probability is $(m-1)/m$ for every label.  Hence
the cap is attained.

For two guesses, first note that
\begin{equation}\label{eq:root-sum}
 \sum_{i<j}\ket{\phi_i-\phi_j}\!\bra{\phi_i-\phi_j}
 =m\sum_i\ket{\phi_i}\!\bra{\phi_i}
   -\ket{\sum_i\phi_i}\!\bra{\sum_i\phi_i}
 =\frac{m^2}{m-1}I.
\end{equation}
Since $\|\phi_i-\phi_j\|^2=2m/(m-1)$, substituting
\eqref{eq:edge-povm} into \eqref{eq:root-sum} gives
$\sum_{i<j}M_{ij}=I$.

If $k\notin\{i,j\}$, then \eqref{eq:simplex} gives
\[
 \langle\phi_i-\phi_j,\phi_k\rangle=0.
\]
Thus an outcome whose reported pair omits $k$ has zero probability on state
$\phi_k$.  Every possible outcome contains the true label, and the two-list
success is one.  No probability can exceed one, completing the proof.
\end{proof}

The pair differences in \eqref{eq:edge-povm} form the normalized root frame of
type $A_{m-1}$.  For $m=3$, the construction is the usual trine exclusion
measurement.  Antidistinguishability and conclusive exclusion are established
subjects~\cite{bandyopadhyay2014}; novelty is not claimed for that isolated
fact.  Its role here is to prove exact attainment of the first nontrivial
union step for an all-distinct family of arbitrary size.

\section{A full union matroid is not sufficient}\label{sec:counterexample}

The condition $r_\ell(E)=M$ says that the support inequalities alone do not
forbid perfect list success.  It does not construct a POVM.  The following
exact example prevents the combinatorial theorem from being misread as a
geometric characterization.

\begin{proposition}[Hemisphere obstruction]\label{prop:counterexample}
Consider the three pure qubit states
\begin{equation}\label{eq:counter-states}
 \phi_1=\begin{pmatrix}1\\0\end{pmatrix},\qquad
 \phi_2=\frac1{\sqrt5}\begin{pmatrix}2\\1\end{pmatrix},\qquad
 \phi_3=\frac1{\sqrt5}\begin{pmatrix}2\\i\end{pmatrix}.
\end{equation}
Their support matroid is $U_{2,3}$, so
$r_{\cR^{\vee2}}(E)=3$ and the two-list support cap equals one.  Nevertheless
no perfect two-list measurement exists.  For uniform priors its exact optimum
is
\begin{equation}\label{eq:counter-opt}
 P_{\rm list}^{\rm opt}=\frac{2+1/\sqrt3}{3}<1.
\end{equation}
\end{proposition}

\begin{proof}
The squared pairwise overlaps are
\begin{equation}\label{eq:counter-overlaps}
 |\langle\phi_1,\phi_2\rangle|^2=\frac45,
 \quad |\langle\phi_1,\phi_3\rangle|^2=\frac45,
 \quad |\langle\phi_2,\phi_3\rangle|^2=\frac{17}{25}.
\end{equation}
Every pair is therefore linearly independent.  The matroid is $U_{2,3}$,
whose two-fold union is the free matroid on three labels.

Suppose a perfect two-list measurement existed.  Every outcome list omits at
least one of the three labels.  Assign each outcome to one omitted label and
sum the corresponding POVM effects, obtaining effects $F_1,F_2,F_3$ with
$\sum_iF_i=I$.  Perfect inclusion implies that every outcome omitting $i$
has zero probability on $\phi_i$, so
\begin{equation}\label{eq:antidist}
 F_i\phi_i=0.
\end{equation}
Thus $(F_i)$ would antidistinguish the three states.

Let $n_i\in\mathbb R^3$ be their Bloch vectors.  A positive qubit operator
annihilating $\phi_i$ has the form
$F_i=a_i\ket{\phi_i^\perp}\!\bra{\phi_i^\perp}$ with $a_i\geq0$.
Completeness requires
\begin{equation}\label{eq:bloch-complete}
 \sum_i a_i=2,\qquad \sum_i a_i n_i=0.
\end{equation}
For \eqref{eq:counter-states}, the $z$ coordinates of the Bloch vectors are
\begin{equation}\label{eq:zpositive}
 (n_1)_z=1,\qquad(n_2)_z=(n_3)_z=\frac35.
\end{equation}
They lie in a strict open hemisphere.  The second equation in
\eqref{eq:bloch-complete} cannot hold with nonnegative $a_i$ whose sum is
two.  This contradicts \eqref{eq:antidist} and proves nonperfectness.

It remains to prove the exact value \eqref{eq:counter-opt}.  For three labels,
minimum list failure equals minimum exclusion error.  An exclusion POVM
$(F_i)$ yields the list $E\setminus\{i\}$ on outcome $i$.  Conversely, from
any two-list POVM assign each outcome to one label omitted by its list; the
resulting exclusion error is no larger than the original list failure.
Taking infima in the two directions proves equality.

The uniform exclusion SDP and its dual are
\begin{equation}\label{eq:exclusion-sdp}
 e_*=\min_{F_i\geq0,\ \sum_iF_i=I}
 \frac13\sum_i\Tr(F_i\rho_i)
 =\max_{Y\leq\rho_i/3}\Tr Y.
\end{equation}
Put
\[
 u=\left(\frac13,\frac13,\frac23\right),\qquad
 R=\frac1{\sqrt3}.
\]
The Bloch vectors are
\[
 n_1=(0,0,1),\quad n_2=(4/5,0,3/5),\quad
 n_3=(0,4/5,3/5),
\]
and satisfy $|u-n_i|=R$ and $u\cdot n_i=2/3$.  The dual operator
\begin{equation}\label{eq:dual-Y}
 Y=\frac16\bigl[(1-R)I+u\cdot\sigma\bigr]
\end{equation}
is feasible because
\[
 \frac13\rho_i-Y
 =\frac16\bigl[R I+(n_i-u)\cdot\sigma\bigr]\geq0.
\]
Its trace is $(1-R)/3$.

For primal attainment, define
\[
 m_i=\frac{u-n_i}{R},\qquad
 (w_1,w_2,w_3)=\left(\frac13,\frac56,\frac56\right),
 \qquad F_i=\frac{w_i}{2}(I+m_i\cdot\sigma).
\]
Here $|m_i|=1$, $\sum_iw_i=2$, and
$\sum_iw_in_i=2u$, hence $\sum_iF_i=I$.  Moreover
$m_i\cdot n_i=-R$, and therefore
\[
 \frac13\sum_i\Tr(F_i\rho_i)
 =\frac13\sum_i\frac{w_i}{2}(1-R)
 =\frac{1-R}{3}.
\]
Primal and dual values coincide.  Thus
$e_*=(1-1/\sqrt3)/3$ and $P_{\rm list}^{\rm opt}=1-e_*$, proving
\eqref{eq:counter-opt}.
\end{proof}

The example also shows why perturbing a degenerate support arrangement can
change attainability without changing the matroid.  Support ranks are
piecewise constant; antidistinguishability is controlled by a convex
geometric condition inside that rank stratum.

\section{Computation and exact certificates}\label{sec:computation}

For a represented subspace family, the cap is finite and combinatorial.  One
may evaluate $d(C)$ by exact matrix rank and use
\eqref{eq:union-rank} directly for small $M$.  For larger instances, standard
matroid-union and submodular-minimization algorithms apply; no quantum SDP is
needed to certify the cap.  Cunningham gives improved algorithms for matroid
partition and intersection~\cite{cunningham1986}.  This computational
observation concerns the upper bound, not construction of an optimal POVM.

The accompanying verifiers use the Python standard library and exact rational
or algebraic arithmetic.  They perform the following fail-closed checks:

\begin{enumerate}
 \item exhaustive comparison of \eqref{eq:union-rank} with a brute-force
       coloring definition on 240 small subspace instances;
 \item all subset inequalities for an exact list-measurement fixture;
 \item the one-guess and edge-root POVMs for regular simplices
       $m=3,\ldots,8$ through rational Gram identities;
 \item exhaustive laminar-channel phase values for $70$ small
       $(h,s,q)$ instances, including the partition-matroid cell formula;
 \item complete-unitary-basis frame spectra, Schmidt-rank caps, Weyl wiring
       values, and the rank-two fixed-probe list certificate;
 \item the exact overlaps, free union matroid, strict-hemisphere obstruction,
       and primal--dual optimum of \Cref{prop:counterexample};
 \item the strict-flat value $81/100$ versus the total-dimension value one;
 \item monotonicity of support-sum ranks under a common process compression.
\end{enumerate}

The finite tests are not a proof of the general theorem.  They guard the rank
formula, normalizations, and fixtures against transcription and convention
errors.  The general proofs are
\Cref{lem:union-rank,thm:main,thm:process,thm:classical-reduction} and
\Cref{thm:laminar-phase,thm:ueb-spectrum,thm:weyl-wiring}.

\section{Related work and priority boundary}\label{sec:related}

Quantum minimum-error discrimination is classically formulated through the
Holevo--Yuen--Kennedy--Lax optimality conditions~\cite{holevo1973,yuen1975}
and is reviewed in standard texts and surveys~\cite{helstrom1976,barnett2009}.
State exclusion reverses the decision: an outcome rules out one or more
hypotheses.  Bandyopadhyay et al. give an SDP, optimality conditions, and
variants including unambiguous and worst-case exclusion
~\cite{bandyopadhyay2014}.  Experimental elimination of multiple states is
demonstrated by Webb et al.~\cite{webb2023}.  Stratton, Hsieh, and Skrzypczyk
derive projector and rank conditions for weak and strong $k$-state exclusion
~\cite{stratton2024}.  Johnston, Russo, and Sikora characterize pure-state
antidistinguishability through Gram-matrix incoherence and give tight overlap
bounds~\cite{johnston2025}.  Recent work studies group-generated exclusion
~\cite{yao2026}, global-versus-LOCC higher-order exclusion~\cite{manna2026},
and zero-error list decoding of classical--quantum channels, including the
trine at list size two~\cite{dalai2026}.  By \Cref{lem:padding}, our list
success is the complementary set-valued formulation with $k=M-\ell$.
Therefore neither list decoding, the simplex/trine measurement, state
exclusion, nor the nonsufficiency of rank conditions is presented as new.

Rado's theorem on independent representatives is classical
~\cite{rado1942}.  Edmonds established the minimum-partition criterion and
matroid union~\cite{edmonds1965,edmonds1968}, as well as the rank-inequality
description of matroid polyhedra~\cite{edmonds1970}.  Our use of these results
is direct and fully credited.  The potentially new point is narrower: the
all-subset Bayesian inequality that places the entire inclusion-probability
vector in an exact union-matroid polytope, together with arbitrary priors,
soft rewards, the robust core-tail version, and canonical process-tester
compression.

Quantum testers and combs provide a common language for discrimination of
multi-time processes~\cite{chiribella2009}.  Canonical inverse-square-root
normalization is standard~\cite{sedlak2016}; general and indefinite-order
strategies have been compared in concrete unitary discrimination problems
~\cite{bavaresco2021,bavaresco2022}, and restricted tester classes have a
general convex formulation~\cite{nakahira2021}.  Entanglement-breaking
channels can already exhibit finite-query adaptive advantages
~\cite{harrow2010}; asymptotic comparisons obey different laws
~\cite{salek2022}.  \Cref{lem:compression} recalls the normalization only to
make the list theorem convention-safe.  The input-dependent application does
not claim adaptivity, binary search, or membership queries as new.  Its
specific contribution is the exact all-depth list phase and the proof that
arbitrary quantum parallel probes and memories reduce to the stated cell and
decision-tree optimizations for this family.

Complete unitary-error bases, dense coding, and symmetric operation
discrimination are likewise established
~\cite{feng2006,sacchi2005,wu2006,zhuang2020}.  In particular, Feng, Duan,
and Ji already derive and attain the approximate dense-coding success law
that becomes \eqref{eq:ueb-success} in our notation.
Independent unitary histories admit strong parallel reductions
~\cite{chiribella2008}.  Accordingly, \Cref{thm:ueb-spectrum,thm:weyl-wiring}
do not claim the twirl, dense-coding protocol, spectrum formula, or
support-dimension converse as new.  The narrower addition is the
process-list translation, its Schmidt-rank equality audit, the list-wiring
trichotomy, and the explicit intermediate-rank failure of the coarse list
cap.

The one-label Rado-matroid support bound was developed separately in
Ref.~\cite{eriksson2026matroid}.  The present paper cites that work rather than
repackaging it.  Its distinct contribution is the operational list hierarchy,
the closed union-rank law, soft rewards, obstruction threshold, exact
attainment/insufficiency pair, and the two architecture-sensitive channel
theorems above.  Targeted searches did not locate this combined theorem
package.  Absence from a search is not proof of priority; no claim of being
first is made.

\section{Limitations and conclusions}\label{sec:limitations}

Several boundaries are essential.

First, the theorem is an upper bound.  The inclusion vector belongs to a
matroid polytope, but not every point of that polytope is generated by a POVM.
\Cref{prop:counterexample} shows that even a free union matroid need not permit
perfect list discrimination.

Second, the theorem uses exact support.  An arbitrarily small full-rank noise
component can enlarge every support and weaken the raw cap discontinuously.
\Cref{cor:robust} mitigates this with chosen positive cores and an explicit
tail penalty, but it does not canonically choose those cores or guarantee that
the resulting certificate improves an SDP bound.

Third, the tester-independent process cap uses the uncompressed Choi support.
A fixed tester normalization can lower the Rado-union rank and give a stronger
certificate.  Conversely, the cap does not determine which causal
architecture attains the actual optimum.  The exact architecture formulas
proved here rely on either complete dephasing with a common input measurement
or a complete unitary-error frame; they do not extend automatically to noisy,
incomplete, or nonuniform channel families.

Fourth, the adaptive channel comparison excludes indefinite causal order,
coherent phase-oracle access, inverse calls, and controlled bypass.  The
$\mathrm{SER}_0$ wiring excludes side memory, intermediate measurements, and
retained transcripts.  These resource definitions are part of the theorems,
not incidental implementation details.

Fifth, the soft-list extension assumes rewards in $[0,1]$ and total reward
budget at most $\ell$ for each outcome.  General cost matrices with negative
entries or larger column sums are outside the theorem.

Within these limits, the result gives a closed answer to a precise question:
what can support dependence alone certify about quantum list decisions?  The
answer is the independence polytope of a matroid union.  It yields arbitrary
prior bounds, an integer congestion deficit, an exact obstruction-disappearance
threshold, and a process-tester version under the stated Choi interface.  The simplex and hemisphere
examples show both sides of the boundary: the cap can be exactly attained by
an all-distinct nonorthogonal family, but support combinatorics alone cannot
replace quantum geometry.  The channel theorems then identify two settings in
which architecture optimization does close exactly: feedback doubles the
number of resolved laminar cells at every call, while complete unitary frames
convert the Schmidt spectrum and wiring dimension into exact success laws.

\section*{Data and code availability}

The accompanying source package contains the \LaTeX{} manuscript, bibliography,
claim ledger, priority audit, research memo, and four standard-library Python
verifiers.  They perform exact rational or algebraic checks and require no
network access or external data.

\section*{AI-assistance disclosure}

The author used AI assistance for literature triage, algebraic cross-checks,
code generation, and editorial revision.  The mathematical claims,
conventions, and final responsibility remain with the author.

\bibliographystyle{unsrt}
\bibliography{references}

\appendix
\section{A direct derivation for deterministic lists}

For completeness, specialize \Cref{thm:main} to a deterministic list POVM.
For $A\subseteq E$,
\begin{align*}
 s(A)
 &=\sum_{i\in A}\sum_{L\ni i}\Tr(M_L\rho_i)\\
 &=\sum_L\Tr\!\left(M_L\sum_{i\in A\cap L}\rho_i\right)\\
 &\leq\sum_L |A\cap L|\Tr(M_LP_A)\\
 &\leq\ell\Tr\!\left(P_A\sum_LM_L\right)
 =\ell d(A).
\end{align*}
For $C\subseteq A$, split $A=C\mathbin{\dot\cup}(A\setminus C)$ and use
$s_i\leq1$ on the second part.  This gives
\[
 s(A)\leq\ell d(C)+|A\setminus C|.
\]
The minimization over $C$ is not an optional relaxation: by
\Cref{lem:union-rank}, it is exactly the rank of $\cR^{\vee\ell}$.

\end{document}
